% Revised from cached-score audit. Regenerate tables with python audit_icassp.py.
% Original manuscript is preserved in manuscript_audit/original/.
\documentclass[10pt,letterpaper]{article}
\usepackage{spconf}
\usepackage{amsmath,amssymb,graphicx,booktabs,cite}
\usepackage{balance}
\usepackage[hidelinks]{hyperref}
\usepackage{url}
\urlstyle{same}
\hypersetup{pdftitle={Class-Prior Constraints on Confident-Tail Test-Time Adaptation for Audio Deepfake Detection},
pdfauthor={Roham Izadidoost, Sumit Srivastava, Shweta Sharma}}
\makeatletter
\newcommand{\tablerows}[1]{\@@input #1 }
\makeatother
\setlength{\textfloatsep}{10pt plus 2pt minus 2pt}
\setlength{\dbltextfloatsep}{10pt plus 2pt minus 2pt}
\title{Class-Prior Constraints on Confident-Tail Test-Time Adaptation\\
for Audio Deepfake Detection}
\name{Roham Izadidoost$^{1}$ \qquad Sumit Srivastava$^{2}$ \qquad Shweta Sharma$^{2}$}
\address{$^{1}$Iran University of Science and Technology \qquad $^{2}$Manipal University Jaipur\\
\small izadidoostroham@gmail.com, sumit.srivastava@jaipur.manipal.edu,\\
\small Shweta.sharma1@jaipur.manipal.edu}
\begin{document}
\maketitle

\begin{abstract}
Confident-tail self-training assigns pseudo-labels to the lowest and highest
scoring target examples. We study a limitation of this rule in audio deepfake
detection: a selected class-$c$ bucket occupying fraction $b_c$ of a pool with
class prior $\pi_c$ has purity at most $\min(1,\pi_c/b_c)$, regardless of the
detector. For exact symmetric tails of size $q$, perfect purity therefore
requires $q\leq\min(\pi,1-\pi)$; this is not a guarantee of successful
adaptation. On 400,435 available ASVspoof2021-DF eval trials, three released
checkpoints fall from 98.78--99.11\% accuracy to 61.80--65.24\% after adaptation
with nominal $q=0.3$. Allocating the tails using a black-box label-shift estimate
restores accuracy to 98.74--98.98\%, while improvements over the source vary by
checkpoint and metric. Controlled prevalence resampling reproduces the large
budget's degradation on In-the-Wild. Smaller tails avoid that large collapse
without estimating a prior, whereas prior estimation itself can fail under
corpus shift. Recomputed tie-aware metrics, threshold-only controls, and disjoint
evaluation distinguish damage prevention from improved discrimination.
\end{abstract}
\begin{keywords}
audio deepfake detection, test-time adaptation, label shift, pseudo-labels
\end{keywords}

\section{Introduction}
\label{sec:intro}
Self-supervised speech representations paired with anti-spoofing back-ends
provide strong audio deepfake detectors~\cite{wav2vec2aug,aasist}, but transfer
across recording conditions and synthesis systems remains difficult~\cite{itw}.
A target-domain failure can involve both score ordering and the decision
threshold. For example, a released detector in our In-the-Wild evaluation has
ROC-AUC 0.9171 but only 39.08\% accuracy at threshold 0.5. This motivates
examining adaptation and re-thresholding separately; it does not establish
that ranking always transfers or that probabilities are calibrated.

Test-time adaptation (TTA) uses unlabeled target observations to update a source
model. Entropy minimization~\cite{tent} and pseudo-label self-training can
reinforce incorrect predictions. We examine a specific self-training rule:
assign the bottom and top score quantiles to bona fide and spoof classes.
With a fixed symmetric budget, the minority bucket can be much larger than the
entire minority population. Even a perfect ranking cannot fill that bucket
with correct pseudo-labels. This obstruction is especially relevant when
benchmark pools and deployment pools have different class proportions.

Class imbalance in adaptation is already established: Label Shift
Adapter~\cite{lsa} estimates target proportions to adapt model parameters, and
RLSbench~\cite{rlsbench} studies adaptation under joint label and conditional
shift. Black-box shift estimation (BBSE)~\cite{bbse} estimates target priors
from a source confusion matrix. Our contribution is narrower: an explicit
counting constraint for confident-tail sampling, an audio checkpoint study of
its consequences, and a comparison of prior-based allocation with smaller tails
and threshold-only controls. T$^2$A~\cite{t2a} instead addresses deepfake TTA
through uncertainty-aware negative learning. We do not claim a new general
label-shift estimator or superiority to specialized label-shift adaptation.

\section{Tail selection and adaptation}
\label{sec:method}
Let $y=1$ denote spoof and $s_\theta(x)\in[0,1]$ the detector's spoof score.
The unlabeled adaptation pool has $N$ clips and unknown class fractions
$\pi_c$. Source labels are available for BBSE; target labels are used only for
evaluation, controlled resampling, and explicitly marked oracle analyses.

\subsection{A necessary condition for pure pseudo-labels}
A bucket $C_c$ assigned pseudo-label $c$, with realized fraction
$b_c=|C_c|/N>0$, contains at most $\pi_c N$ correct examples. Thus
\begin{equation}
 \operatorname{purity}(C_c)\leq\min\{1,\pi_c/b_c\}.
 \label{eq:bound}
\end{equation}
This is a finite-pool counting bound and makes no assumption about the model.
For exact symmetric tails, $b_0=b_1=q$, perfect purity is possible only if
\begin{equation}
 q\leq\min\{\pi,1-\pi\},\qquad \pi=\pi_1.
 \label{eq:necessary}
\end{equation}
For illustration, at 97.22\% spoof prevalence a bucket labeling 30\% of the pool
bona fide has purity at most 9.27\%. High AUC cannot remove this constraint.
Satisfying it only removes the counting obstruction: overlapping score
distributions can still contaminate both tails, and pure tails do not guarantee
beneficial gradient updates. Conversely, learning may tolerate impure labels.

A prior allocation uses nominal fractions $q_c=2q\hat\pi_c$, keeping the
nominal total selection fraction $2q$. With exact bucket sizes and a correct
prior, $q\leq1/2$ ensures $q_c\leq\pi_c$. With estimated priors, the
corresponding feasibility condition is
\begin{equation}
 \hat\pi_c/\pi_c\leq1/(2q)\quad(c=0,1;\ \pi_c>0).
 \label{eq:tolerance}
\end{equation}
It limits permissible overestimation of each class; it is not a safety theorem.
A smaller symmetric $q$ is an alternative requiring no point estimate of the
prior, but choosing a guaranteed feasible nonzero value still requires a lower
bound on the minority prevalence.

\textbf{Quantile ties.} Our cached experiments use inclusive score thresholds,
$s\leq Q(q_0)$ and $s\geq Q(1-q_1)$, with spoof assignments taking precedence
if masks overlap. Consequently, realized sizes can differ from nominal budgets.
Equation~\eqref{eq:bound} applies to the realized $b_c$; equations
\eqref{eq:necessary}--\eqref{eq:tolerance} describe exact-size tails. We audit
realized selections on cached source scores and do not interpret nominal
budgets as exact sample counts. Deterministic selection of fixed counts would
be a different implementation requiring new adaptation runs.

\subsection{Prior estimation and objective}
For $M_{ij}=P_{\mathcal S}(\hat y=j\mid y=i)$ and target prediction histogram
$h$, BBSE solves $M^\top\hat p=h$. We estimate $M$ using up to 4,000
class-balanced clips from the checkpoint's labeled training corpus; $h$ uses
source predictions on the adaptation pool. We clip the solution components to
$[10^{-4},1-10^{-4}]$ and renormalize. The implementation returns a uniform
prior if the matrix is singular. Because $M$ uses training rather than held-out
source data, optimism in its estimated error rates is a limitation.

At each of four passes, the current model scores the adaptation pool and
refreshes pseudo-labels. For a mini-batch $B$ and its selected subset $B_C$,
\begin{equation}
 \begin{aligned}
 \mathcal L_B ={}&
 \frac{1}{|B_C|}\sum_{x\in B_C}\!\mathrm{CE}(p_\theta(x),\tilde y_x)
 \\ &+\frac{\lambda}{2|B|}\sum_{x\in B}
 \|p_\theta(a(x))-\mathrm{sg}[p_\theta(x)]\|_2^2,
 \end{aligned}
 \label{eq:loss}
\end{equation}
where $p_\theta=(1-s_\theta,s_\theta)$, the CE term is omitted when $B_C$ is
empty, and $\lambda=0.3$. The consistency term uses all batch examples and both
probability components. The perturbation is per-clip gain sampled uniformly
from $[0.7,1.3]$ plus Gaussian noise with standard deviation 0.005. Only the
last four transformer blocks' LayerNorm affine parameters and the classifier
head are updated; BatchNorm statistics stay frozen. Adam uses learning rate
$10^{-4}$ and batch size 32. The released DeepFense models have 16,706 trainable
parameters under this rule. A fixed-confidence control instead selects scores
at most 0.05 or at least 0.95, with the same objective.

\section{Evaluation protocol}
\label{sec:protocol}
\textbf{Checkpoints and pools.} We use DeepFense XLS-R 300M + AASIST
checkpoints trained on ASVspoof2019-LA with training seeds 2, 42, and
240~\cite{deepfense,deepfenseweights}, denoted DF-2, DF-42, and DF-240. A
WaveFake-trained SSL-AASIST checkpoint, WF~\cite{shiftyspeech}, is evaluated
on In-the-Wild. Audio is 16\,kHz with the checkpoint-specific input length:
64,000 samples for DeepFense and 64,600 for WF.

The official ASVspoof2021-DF metadata~\cite{asvspoof21} contains 611,829 trials
across three phases, comprising 533,928 eval trials, 59,325 progress trials,
and 18,576 hidden trials. Our
available \emph{eval-only} subset contains 400,435 trials (75.00\% coverage):
389,299 spoof and 11,136 bona fide, or 97.219\% spoof. Coverage by class is
75.00\% spoof and 74.89\% bona fide. The missing shard prevents a full-set
comparison; these proportions do not establish an unbiased estimate of full-set
EER. All method comparisons use the same available trials.

The main run adapts on 20,000 randomly selected clips with adaptation seed 0.
We report both the available pool, which includes adaptation clips, and its
380,435-clip disjoint remainder. Additional adaptation seeds 1 and 2 use a fixed
50,000-trial subset, selected with seed 12345, with 20,000 clips for adaptation
per run. These are sensitivity checks on overlapping data, not nine independent
datasets. In-the-Wild uses all 31,779 clips (11,816 spoof; 37.182\%), with
20,000 adaptation clips and an 11,779-clip disjoint remainder.

\textbf{Prevalence intervention.} On In-the-Wild with WF, we retain all spoof
clips and randomly thin bona fide clips to spoof proportions 0.90, 0.95, and
0.97 (sampling seed 4242). These pools are smaller than 20,000, so adaptation
and evaluation use the entire resampled pool. The experiment changes both
prevalence and pool size; it is transductive and does not isolate prevalence
from sample-count effects.

\begin{table*}[t]
\centering
\caption{Recomputed results on the 400,435 available DF eval trials, checkpoint
DF-2, adaptation seed 0. Acc. and balanced accuracy (BA) use threshold 0.5 except
for the two frozen-source threshold controls. $G_\tau$ is an oracle diagnostic
in percentage points; smaller is better. Baselines are the specified variants
in Sec.~\ref{sec:protocol}.}
\label{tab:df}
\small
\begin{tabular}{@{}lrrrrr@{}}
\toprule
Method & EER (\%) & AUC & Acc. (\%) & BA (\%) & $G_\tau$\\
\midrule
\tablerows{manuscript_audit/table_df_main.tex}
\bottomrule
\end{tabular}
\end{table*}

\textbf{Metrics and controls.} We recompute metrics from saved float32
probabilities. ROC-AUC assigns half credit to ties; EER linearly interpolates
the crossing of FPR and FNR on the ROC, grouping identical scores. Interpolated
EER need not correspond to an attainable deterministic threshold. We also
report accuracy, balanced accuracy, and
\begin{equation}
 G_\tau=\max_t\operatorname{Acc}(s,t)-\operatorname{Acc}(s,\tau).
 \label{eq:gap}
\end{equation}
The maximum considers all distinct score thresholds, including constant
predictions. It uses evaluation labels and is strictly a diagnostic of
threshold headroom, not a deployable method or a probability-calibration error.
Unlike $1-\mathrm{EER}$, it is the actual attainable accuracy maximum on the
observed scores.

Two label-free controls keep source scores fixed and select a threshold from
adaptation scores: their median, or their $(1-\hat\pi)$-quantile using BBSE.
The latter is a predicted-prevalence matching rule, not a guarantee of Bayes
optimality. These controls preserve AUC and EER. Their cached-score prior
estimates can differ slightly from the original adaptation-time forward passes.

We additionally report fixed-configuration entropy adaptation (Tent-style), an
ETA reliability-only ablation without redundancy or Fisher regularization,
and SAR with recovery threshold scaled to $(0.2/\ln1000)\ln2$~\cite{tent,eata,sar}.
These variants share the above update settings. IM-PL is a SHOT-inspired
information-maximization control with probability-space centroid pseudo-labels,
not the original feature-space SHOT algorithm~\cite{shot}; it freezes the head
and updates 16,384 parameters. Thus these are controlled variants, not tuned
reproductions establishing superiority over the published methods.

\section{Results}
\label{sec:results}
\subsection{Preventing large-budget damage on DF}
Table~\ref{tab:df} shows that symmetric $q=0.3$ lowers DF-2 accuracy from
98.87\% to 62.96\% and increases EER from 4.51\% to 5.84\%. Its oracle
threshold gap is 35.38 points: much of the fixed-threshold damage is
recoverable by re-thresholding, although AUC also declines. The BBSE budget
reaches 98.95\% accuracy and 4.06\% EER. Fixed confidence reaches higher
accuracy (99.02\%), and the smaller $q=0.02$ reaches 98.96\% accuracy and
4.19\% EER. Prior allocation is therefore one effective remedy, not a uniformly
best choice. Balanced accuracy is lower after BBSE adaptation than for the
source (82.34\% versus 88.28\%); the slight overall accuracy gain should not
be interpreted as an improvement for both classes.

\begin{table}[t]
\centering
\caption{DF checkpoint comparison on the same available eval pool. Entries
are EER\%/accuracy\%; symmetric and BBSE budgets use nominal $q=0.3$.}
\label{tab:dfseeds}
\small\setlength{\tabcolsep}{3pt}
\begin{tabular}{@{}lccc@{}}
\toprule
Checkpoint & Source & Sym. & BBSE\\
\midrule
\tablerows{manuscript_audit/table_df2021_checkpoints.tex}
\bottomrule
\end{tabular}
\end{table}

The accuracy collapse and its reversal occur on all three checkpoints
(Table~\ref{tab:dfseeds}). Across the nine checkpoint/run combinations, BBSE
has lower EER than the symmetric budget in 9/9 and the source in 8/9. We give
these as descriptive counts without treating reused checkpoints and trials as
independent evidence. The disjoint evaluation is more qualified: BBSE changes
source EER from 4.52 to 4.07 (DF-2), 4.54 to 3.29 (DF-42), and 3.83 to 3.83
(DF-240, a 0.003-point increase before rounding). Avoiding severe damage is
more consistent than improving an already strong source.

The frozen-source BBSE threshold gives 98.88\% accuracy for DF-2, close to its
98.87\% source result, while the median gives 59.06\%. Median thresholding
implicitly targets roughly balanced predictions and is poorly suited to this
pool. The budget change can improve EER, which threshold changes cannot do,
but it is not required to preserve high accuracy here. Removing consistency
from the BBSE objective gives 3.98\% EER versus 4.06\% with it on DF-2; this
single ablation provides no evidence that consistency is needed for the gain.

\subsection{Score resolution and baseline interpretation}
Tent-style entropy adaptation produces 98.07\% of scores within $10^{-6}$ of
1.0, but only 93.02\% exactly equal 1.0. ETA's corresponding fractions are
97.71\% and 96.65\%. Conflating these quantities would make the stated AUCs
appear impossible. With exact ties grouped, Tent-style adaptation has EER
4.28\% and AUC 0.9765, compared with source 4.51\% and 0.9923: a small
improvement near one ROC operating point can coexist with worse integrated
ranking. For ETA, ROC interpolation gives EER 12.31\%; selecting the nearest
ROC point and averaging its errors would give 7.44\%, illustrating sensitivity
to the scoring convention. For DF-42 and DF-240 its balanced accuracy
is 50\%, indicating an all-spoof decision at threshold 0.5. These observations
support reporting several metrics, not dismissing EER as arbitrary tie-breaking.

\subsection{Corpus shift and prior-estimation failure}
\begin{table}[t]
\centering
\caption{Native In-the-Wild, EER\%/accuracy\%, with nominal $q=0.3$. WF is the
WaveFake-trained checkpoint; all other rows use ASVspoof2019-trained models.}
\label{tab:itw}
\small\setlength{\tabcolsep}{3pt}
\begin{tabular}{@{}lccc@{}}
\toprule
Checkpoint & Source & Sym. & BBSE\\
\midrule
\tablerows{manuscript_audit/table_itw_checkpoints.tex}
\bottomrule
\end{tabular}
\end{table}
At the native In-the-Wild prior, the nominal symmetric budget has no counting
obstruction. It improves accuracy on three checkpoints but reduces DF-2
accuracy (Table~\ref{tab:itw}), consistent with feasibility being insufficient
for a benefit. For DF-42, symmetric adaptation raises AUC from 0.9171 to
0.9473 and accuracy from 39.08\% to 82.39\%. A frozen-source median threshold
already reaches 78.10\%, showing how much the comparison depends on the
operating point. The oracle threshold gap falls from 48.28 to 8.34 points.

BBSE instead estimates approximately 97.9\% spoof on this 37.2\%-spoof corpus
and obtains 14.54\% EER and 47.99\% accuracy. A two-component score-mixture
estimate gives 14.91\%/55.00\%. An oracle budget using the adaptation pool's
true prior gives 11.59\%/91.28\%. This isolates a substantial estimation cost
on this checkpoint; the oracle is not a deployable method and is not always
best on other checkpoints (WF: 4.67\% oracle EER versus 4.55\% symmetric).

BBSE is identifiable under label shift with an invertible confusion matrix,
even for uncalibrated predictors~\cite{bbse}. Under unrestricted conditional
shift, the hard prediction rate instead satisfies
\begin{equation}
 P_{\mathcal T}(\hat y=1)=\pi\,\mathrm{TPR}_{\mathcal T}
 +(1-\pi)\,\mathrm{FPR}_{\mathcal T},
\end{equation}
which alone cannot determine all three unknowns. This explains why a source
confusion matrix may fail to transfer. It does not prove impossibility for
all output-based estimators with additional assumptions or side information,
nor equivalence between prior error and calibration error.

\subsection{Intervening on prevalence}
\begin{table}[t]
\centering
\caption{WF on resampled In-the-Wild: accuracy\% for the unchanged source and
symmetric-tail adaptation. Pool size varies with prevalence.}
\label{tab:skew}
\small\setlength{\tabcolsep}{3pt}
\begin{tabular}{@{}rrrrrr@{}}
\toprule
$\pi$ & $N$ & Source & $q=.02$ & $q=.10$ & $q=.30$\\
\midrule
\tablerows{manuscript_audit/table_skew.tex}
\bottomrule
\end{tabular}
\end{table}
The large-budget degradation also occurs on resampled In-the-Wild
(Table~\ref{tab:skew}). At 97\% spoof, nominal $q=0.3$ permits at most 10\%
purity in an exact-size bona fide tail and yields 65.42\% accuracy; $q=0.02$
yields 98.59\%. However, the unchanged source reaches 98.94\% there, so
small $q$ prevents the large collapse without universally improving the source.
The experiment supports prevalence sensitivity of this selection rule; it
neither proves inevitable failure whenever the bound is below one nor identifies
prevalence as the only cause of adaptation failure.

\section{Discussion and limitations}
The counting obstruction gives a useful diagnostic for pseudo-label allocation,
while its consequences depend on score quality, realized selection sizes, and
the learning objective. Recomputed metrics support severe fixed-threshold
damage from large symmetric tails and its mitigation by prior allocation or
smaller tails. They also show important tradeoffs: some remedies reduce
minority-class performance, prior estimates can fail on shifted corpora, and
re-thresholding alone explains a substantial part of some accuracy gains.

The experiments use four released checkpoints from one broad architecture
family, two principal corpora, an incomplete DF eval partition, and fixed
adaptation hyperparameters. They do not establish general superiority over
specialized label-shift methods. Our prevalence intervention changes sample
count, and the seed checks share data. Future validation should use independent
corpora, fixed-size prevalence interventions, held-out source calibration,
exact-count tail selection, and development-domain tuning of baseline methods.
Neither a small nominal tail nor source-score agreement alone certifies safe
adaptation without assumptions about the target distribution.

\textbf{Reproducibility.} The script \texttt{audit\_icassp.py} verifies 89 saved score arrays
against historical AUC and accuracy, recomputes the reported metrics, and
writes the tables and input hashes. The project repository is
\url{https://github.com/RohamIzadidoost/Kaggle_deepFakeAudio}.
The accompanying audit bundle includes the revised scripts and outputs.
No new neural-network training is required
to reproduce these tables from the saved scores and original manifests.

\textbf{Ethics and assistance.} This study uses existing public datasets and
checkpoints; no new human-subject data were collected. OpenAI Codex assisted
with the revised manuscript text, metric-audit code, and result tables; the
adaptation experiments predate this revision.

\clearpage
\balance
\begin{thebibliography}{99}
\bibitem{wav2vec2aug} H.~Tak \emph{et al.}, ``Automatic speaker verification
spoofing and deepfake detection using wav2vec 2.0 and data augmentation,'' in
\emph{Proc. Odyssey}, 2022.
\bibitem{aasist} J.~Jung \emph{et al.}, ``AASIST: Audio anti-spoofing using
integrated spectro-temporal graph attention networks,'' in \emph{Proc. ICASSP}, 2022.
\bibitem{itw} N.~M\"uller, P.~Czempin, F.~Dieckmann, A.~Froghyar, and
K.~B\"ottinger, ``Does audio deepfake detection generalize?,'' in
\emph{Proc. Interspeech}, 2022.
\bibitem{tent} D.~Wang \emph{et al.}, ``Tent: Fully test-time adaptation by
entropy minimization,'' in \emph{Proc. ICLR}, 2021.
\bibitem{lsa} S.~Park, S.~Yang, J.~Choo, and S.~Yun, ``Label shift adapter for
test-time adaptation under covariate and label shifts,'' in \emph{Proc. ICCV},
2023. \url{https://arxiv.org/abs/2308.08810}.
\bibitem{rlsbench} S.~Garg, N.~Erickson, J.~Sharpnack, A.~Smola,
S.~Balakrishnan, and Z.~C.~Lipton, ``RLSbench: Domain adaptation under relaxed
label shift,'' in \emph{Proc. ICML}, vol.~202, pp.~10879--10928, 2023.
\bibitem{bbse} Z.~Lipton, Y.-X.~Wang, and A.~Smola, ``Detecting and correcting
for label shift with black box predictors,'' in \emph{Proc. ICML}, vol.~80,
pp.~3122--3130, 2018.
\bibitem{t2a} H.-H.~Nguyen-Le, V.-T.~Tran, D.-T.~Nguyen, and N.-A.~Le-Khac,
``Think twice before adaptation: Improving adaptability of deepfake detection
via online test-time adaptation,'' in \emph{Proc. IJCAI}, pp.~7686--7694, 2025.
\bibitem{deepfense} Y.~El~Kheir \emph{et al.}, ``DeepFense: A unified, modular,
and extensible framework for robust deepfake audio detection,''
\emph{arXiv:2604.08450}, 2026.
\bibitem{deepfenseweights} DeepFense, ``ASV19\_\allowbreak Wav2Vec2\_\allowbreak AASIST\_\allowbreak NoAug\_\allowbreak Seed\{2,42,240\},''
released model weights, Hugging Face. \url{https://huggingface.co/DeepFense}.
\bibitem{shiftyspeech} A.~Garg \emph{et al.}, ``ShiftySpeech: A large-scale
synthetic speech dataset with distribution shifts,'' \emph{arXiv:2502.05674},
2025. Released checkpoint: \url{https://huggingface.co/ash56/ssl-aasist}.
\bibitem{asvspoof21} J.~Yamagishi \emph{et al.}, ``ASVspoof 2021: Accelerating
progress in spoofed and deepfake speech detection,'' in \emph{Proc. ASVspoof
2021 Workshop}, pp.~47--54, 2021.
\bibitem{eata} S.~Niu \emph{et al.}, ``Efficient test-time model adaptation
without forgetting,'' in \emph{Proc. ICML}, 2022.
\bibitem{sar} S.~Niu \emph{et al.}, ``Towards stable test-time adaptation in
dynamic wild world,'' in \emph{Proc. ICLR}, 2023.
\bibitem{shot} J.~Liang, D.~Hu, and J.~Feng, ``Do we really need to access the
source data? Source hypothesis transfer for unsupervised domain adaptation,''
in \emph{Proc. ICML}, 2020.
\end{thebibliography}
\end{document}
